% EE6221 CA2 -- state-of-the-art review: humanoid robots (DRAFT)
% Compile: xelatex ca2-humanoid-robots-report.tex (twice, for references and ToC)
\documentclass[11pt,a4paper]{article}
\usepackage[margin=2.3cm]{geometry}
\usepackage{amsmath}
\usepackage{booktabs}
\usepackage{tabularx}
\usepackage{array}
\usepackage{enumitem}
\usepackage{float}
\usepackage{xcolor}
\usepackage{url}
\expandafter\def\expandafter\UrlBreaks\expandafter{\UrlBreaks\do\-}
\usepackage[colorlinks=true,linkcolor=black,citecolor=blue!50!black,urlcolor=blue!50!black]{hyperref}
\setlist{itemsep=2pt,topsep=4pt}
\setlength{\parskip}{4pt}
% Red bracketed text marks passages the author must rewrite in their own voice.
\newcommand{\yours}[1]{\textcolor{red}{[#1]}}

\title{\textbf{EE6221 CA2 Assignment}\\[4pt]
\large State of the Art of General-Purpose Humanoid Robots:\\
Robot Brains and the Limits of Real-World Deployment}
\author{Han Chang \quad G2602636F}
\date{November 2026}

\begin{document}
\maketitle

\begin{abstract}
\noindent
This report reviews the current state of general-purpose humanoid robots,
concentrating on the control stack that turns camera images and a spoken
instruction into coordinated motion of every joint.  I compare commercial
platforms (Boston Dynamics Atlas, Figure~03, Apptronik Apollo, Unitree~G1),
recent research on learned whole-body control and vision-language-action (VLA)
models (GR00T~N1, $\pi_0$, Gemini Robotics 1.5, Helix~02, FAST, HuGe), and the
newer practice of using a general-purpose frontier model as the robot's brain,
mainly OpenAI's GPT-6 Astra but also Claude, Gemini and Grok.  For each I ask
where the model runs, on the robot or in the cloud, and how fast it reacts,
because many applications need offline, on-device operation inside a fixed
latency budget.  In my view the best current solution is a \emph{hierarchical,
learned} architecture in which a slow reasoning model, a fast visuomotor policy
and a very fast whole-body controller are trained to work together, as in
Figure's Helix~02, with Atlas the most deployment-ready hardware.  I also
discuss why humanoids are still mostly at the pilot stage while AI agents are
already used widely in software, the main open challenges (data, reliability,
safety certification, energy and evaluation), likely future directions and what
I learned.
\end{abstract}

\setcounter{tocdepth}{1}
\tableofcontents
\newpage

%-----------------------------------------------------------------
\section{Introduction and scope}

A humanoid robot has a human-like body: two legs, a torso, two arms and, in
most current designs, multi-fingered hands.  The case for the form is
practical.  Factories, warehouses and homes were built for human bodies, so a
robot with the same reach, height and hand shape can in principle use existing
tools, shelves and doorways without redesigning the workplace.  The same line
of argument explains why humanoids are hard: a body with 50 or more actuated
joints that must balance on two small feet is underactuated, has highly coupled
dynamics, and must keep doing so while the hands interact with unknown objects.

For many years the hardware was the bottleneck.  Today several companies can
build reliable humanoid bodies, and the limiting factor has moved to the
``brain'': how the robot perceives a scene, decides what to do, and converts
that decision into stable whole-body motion.  I review the overall humanoid
solution here, but the control stack is the part I examine in detail.  It
connects directly to EE6221: forward and inverse kinematics of redundant arms,
robot dynamics and control, vision sensors and pose estimation, sensor fusion
with the Kalman filter, and vision-based control all appear as layers inside a
modern humanoid.

I compared systems using public technical reports, peer-reviewed or preprint
papers, and official company announcements.  Company announcements are treated
as \emph{claims}, and I note where a result lacks independent verification.
Each solution is judged on five criteria: breadth of tasks it can perform
without task-specific programming; length and reliability of autonomous
operation; data efficiency and ability to transfer across robots; readiness for
real deployment (hardware maturity, safety, integration); and openness and
quality of the evidence.

%-----------------------------------------------------------------
\section{Background: the layers of a humanoid control stack}

Any humanoid is a stack of loops running at very different rates, from the
joint loop at the bottom to the task planner at the top.

\begin{enumerate}[label=\textbf{L\arabic*}]
  \item \textbf{Actuation and state estimation} (about 1\,kHz).  Joint torque
  or current control, plus estimation of the floating base pose from IMU,
  joint encoders and foot contact, typically with a Kalman-type filter.  The
  base pose is not measured directly, so this estimate is what every higher
  layer trusts.
  \item \textbf{Whole-body control (WBC)} (about 50--500\,Hz).  Given a
  desired motion of the hands, feet or centre of mass, compute joint commands
  that achieve it while keeping balance and respecting joint and contact
  limits.
  \item \textbf{Visuomotor policy} (about 10--200\,Hz).  Map camera images and
  proprioception to short-horizon targets such as hand poses, grasp
  commands, or a reference motion for the WBC.
  \item \textbf{Task reasoning} (about 1--10\,Hz or slower).  Interpret a
  language instruction, break it into steps, track progress, and recover from
  failure.
\end{enumerate}

\noindent
Classical humanoids built L2 from models: zero-moment-point (ZMP) walking,
model predictive control (MPC) of a simplified centroidal model, and a
quadratic-program WBC that resolves the redundancy of the kinematic chain
(the same null-space idea used for redundant manipulators, but with contact
and balance constraints added).  L3 and L4 were hand-written state machines.
Over the last two years every layer above L1 has been replaced or augmented by
learned models, which makes the way those layers are divided and connected the
design question that matters most.

%-----------------------------------------------------------------
\section{Review of the state of the art}

\subsection{Hardware platforms}

Boston Dynamics' Atlas (product version, January 2026) is the most
deployment-ready body: 56 degrees of freedom, 50\,kg payload, water
resistance, operation from $-20$ to $40\,^\circ$C, autonomous battery swapping,
and connection to factory software through the Orbit fleet platform.  All
2026 production goes to Hyundai and Google DeepMind~\cite{bd-atlas-2026}.
Figure~03 (October 2025) was redesigned for volume manufacturing and homes,
with palm cameras and tactile hands, and is the body for
Helix~02~\cite{figure-03}.  Apptronik's Apollo is one of the robots used by
Gemini Robotics~1.5~\cite{gemini-15}, and the low-cost Unitree~G1 has become
the common research platform~\cite{huge-2026,motionvla-2026}.

\subsection{Learned whole-body control}

Reinforcement learning in GPU-parallel simulation, followed by sim-to-real
transfer, is now the standard way to obtain robust humanoid locomotion.  A
newer line of work is the \emph{behaviour foundation model}, a whole-body
controller pre-trained on large motion datasets so it can be reused across
tasks~\cite{bfm-survey}.  FAST~\cite{fast-2026} adds a small, constrained
residual policy for new motions and centre-of-mass objectives, which shows
that learned controllers still benefit from classical balance quantities.

\subsection{Vision-language-action models}\label{sec:vla}

A VLA model starts from a vision-language model pre-trained on internet data
and is trained further to output robot actions.  $\pi_0$ (Physical
Intelligence, 2024)~\cite{pi0} is a vision-language backbone with a
flow-matching action expert, trained across many robot types, and it set the
recipe later models follow.  GR00T~N1 (NVIDIA, 2025)~\cite{groot-n1} is an
open-weight humanoid model with a dual-system design: a vision-language
``System~2'' and a diffusion action ``System~1''.  Gemini Robotics~1.5 (Google
DeepMind, 2025)~\cite{gemini-15} pairs a VLA with a separate embodied-reasoning
planner and adds motion transfer, which runs the same skill on different
robots, including Apollo.  Helix~02 (Figure, 2026)~\cite{helix-02} adds a
learned whole-body controller (``System~0'') beneath Helix's System~1 and
System~2; Figure showed a four-minute autonomous dishwasher task across a
kitchen and later claimed multi-hour shifts, and both are company claims
without independent evaluation.  HuGe (October 2026)~\cite{huge-2026}
pre-trains humanoid VLAs on over 2{,}000\,h of egocentric human and robot video
and outputs motion tokens for a pre-trained whole-body controller; on a real G1
the human data improved pick-and-drop success.

\subsection{General-purpose frontier models as the robot brain}\label{sec:astra}

A newer idea, which appeared in late 2026, skips robot-specific training
entirely.  A general-purpose frontier model, the same kind of model behind
today's chatbots and AI agents, is connected to a robot and asked to control
it.  The most tested example is OpenAI's GPT-6 Astra, which several research
groups studied within a few weeks of each other.

HomeBody (Stanford and Caltech, September 2026)~\cite{homebody-2026} runs a
Unitree G1 with Astra.  Astra does not output joint commands: it picks a skill
(grasp, navigate, open a drawer) and its target through a structured tool call,
and the skills run on a local laptop GPU.  The robot tidied an unfamiliar
kitchen, but the authors note pauses between skills caused by Astra's reasoning
latency, extra API cost and overheating finger servos.  It is a project page
with code, not a paper with measured results.

Astranomy (NVIDIA, October 2026)~\cite{astranomy-2026} tests Astra as a robot
policy in simulation without any robot training.  Used as an agent with tools,
it reached a 74.8\% success rate on the RoboLab-120 benchmark, against 28.0\%
for $\pi_0$'s successor $\pi_{0.5}$.  The cost is speed and money: a single
direct call takes about 18.6~s on average, and the agent version uses about
30.9 million tokens (about US\$36) per episode.

SenseNova's joint study (October 2026)~\cite{sensenova-2026} evaluates Astra on
8 robot configurations and 12 benchmarks, including humanoids.  It finds that
the best setup uses Astra as a slow ``System~2'' planner coordinating a fast
``System~1'' controller such as $\pi_{0.5}$, which it calls ``costly but
effective''.  The authors conclude that Astra is not yet practical embodied
intelligence, because of cost, execution time and weak recovery from
irreversible failures.

A general model can therefore understand scenes and plan tasks surprisingly
well with no robot training.  It is far too slow, though, to close a balance or
grasping loop, so in every humanoid example it sits at the top of the stack and
calls faster controllers below it.  All three sources are preprints or project
pages, and I did not find an official technical report from OpenAI describing
Astra's design.

Other general models are being tried as a robot brain in the same way.  Google
builds its own robotics models on top of Gemini: Gemini Robotics-ER~1.5 is a
Gemini model tuned for embodied reasoning and used as the planner above the
Gemini Robotics VLA (Section~\ref{sec:vla})~\cite{gemini-15}.  Robocurve
compared Claude Fable~5, Fable~5.1 and Astra on the same pair of 6-DoF YAM
arms, with the model sending end-effector poses
directly~\cite{robocurve-astra}.  On placing a block in a bowl (20 trials
each), Astra succeeded 19 times, Fable~5.1 8 times and Fable~5 once; on a
harder puzzle task, Astra and Fable~5.1 each succeeded only 2 times in 20.
The authors note the runs were not interleaved, used different rigs, and were
graded by an operator, so this is a small, early comparison.  Tesla plans to
use xAI's Grok as the conversational ``brain'' of its Optimus humanoid, while a
vision network adapted from Tesla's self-driving software handles perception
and movement~\cite{grok-optimus,investorplace-optimus}.  Musk has said that for
a Grok-level conversation Optimus needs a Wi-Fi or cellular connection to the
cloud~\cite{optimus-wiki}, which is the on-device versus cloud trade-off of
Section~\ref{sec:deploy}.  That is a plan, not a published or tested system.

In every case the general model supplies language understanding and planning,
while something faster and more specialised moves the body.

\subsection{Comparison}

Table~\ref{tab:compare} summarises the solutions against the criteria in the
introduction.


\begin{table}[H]
\centering\small
\renewcommand{\arraystretch}{1.25}
\begin{tabularx}{\textwidth}{>{\raggedright\arraybackslash}p{2.3cm}
  >{\raggedright\arraybackslash}X >{\raggedright\arraybackslash}X
  >{\raggedright\arraybackslash}p{2.6cm}}
\toprule
Solution & Architecture & Main strength & Evidence \\
\midrule
Atlas (product) & Industrial humanoid; DeepMind models being integrated &
Deployment readiness: payload, environment rating, battery self-swap, fleet
software & Company specs; first fleets in 2026 \\
GR00T N1 & Dual-system VLA (VLM + diffusion action head) & Open weights;
real + human-video + synthetic data & Paper; sim and real-robot tests \\
Gemini Robotics 1.5 & VLA + separate embodied-reasoning model & Reasoning
before acting; cross-robot motion transfer & Technical report \\
Helix 02 + Figure 03 & Three tiers: S2 reasoning, S1 visuomotor, S0 learned
WBC & Single learned system from pixels to whole-body action; long-horizon
demo & Company demos only \\
FAST & Pre-trained WBC + constrained residual adaptation & Fast adaptation
without forgetting; CoM-aware balance & Preprint; sim and real \\
HuGe & VLA $\rightarrow$ latent motion tokens $\rightarrow$ WBC & Scales data
with egocentric human video & Preprint; benchmark + G1 \\
General frontier models (GPT-6 Astra, Claude Fable, Gemini-ER, Grok) &
General model as planner, calling skills or a fast policy & No robot
training needed; strong task planning; but slow and costly & Preprints,
project pages, small lab tests; Grok only announced \\
\bottomrule
\end{tabularx}
\caption{Comparison of the reviewed humanoid solutions.}
\label{tab:compare}
\end{table}

%-----------------------------------------------------------------
\section{Deployment constraints: where the brain runs and how fast it must react}\label{sec:deploy}

Two practical questions decide whether a solution from Section~3 can be used at
all.  First, what is the \emph{latency budget}, the longest allowed delay
between a change in the environment and the robot's response?  Second, must
the brain run \emph{on-device}, on the robot's own computer, and must it work
\emph{offline}, with no network connection?  The answers depend on the
application, and they rule out some models entirely.

Table~\ref{tab:budget} gives my estimate of the budget for the situations that
matter most.


\begin{table}[H]
\centering\small
\renewcommand{\arraystretch}{1.25}
\begin{tabularx}{\textwidth}{>{\raggedright\arraybackslash}p{3.0cm}
  >{\raggedright\arraybackslash}X >{\raggedright\arraybackslash}p{2.4cm}
  >{\raggedright\arraybackslash}p{3.6cm}}
\toprule
Situation & Example & Latency budget & Where the brain can run \\
\midrule
Balance, fall recovery & Slipping or being pushed & about 1--10\,ms &
On-device only \\
Contact and grasp correction & An object slips in the hand & about
10--50\,ms & On-device \\
Reacting to people & A person steps into the robot's path & about
100\,ms & On-device \\
Task planning & ``Tidy the kitchen'' & about 1--10\,s & Cloud is acceptable
if a network is available \\
Offline sites & Factory with no internet or strict data privacy; disaster
site & Whole stack & Everything on-device, so smaller models \\
\bottomrule
\end{tabularx}
\caption{Indicative latency budgets by situation.  These are my estimates,
based on the loop rates in Section~2, not values from a single source.}
\label{tab:budget}
\end{table}

Table~\ref{tab:sizes} lists what each source states about model size, speed
and where the model runs; many companies do not publish these numbers.


\begin{table}[H]
\centering\small
\renewcommand{\arraystretch}{1.25}
\begin{tabularx}{\textwidth}{>{\raggedright\arraybackslash}p{2.3cm}
  >{\raggedright\arraybackslash}p{2.3cm} >{\raggedright\arraybackslash}X
  >{\raggedright\arraybackslash}p{2.7cm} >{\raggedright\arraybackslash}p{2.4cm}}
\toprule
System & Parameters & Rate / delay & Where it runs & Training data \\
\midrule
Helix~\cite{helix-2025} & S2: 7B; S1: 80M & S2: 7--9\,Hz; S1: 200\,Hz &
Fully onboard, embedded low-power GPUs & About 500\,h teleoperation \\
Helix 02~\cite{helix-02} & S0: 10M; others not disclosed & S1: 200\,Hz;
S0: 1\,kHz & Not stated & Over 1{,}000\,h human motion (for S0) \\
GR00T N1~\cite{groot-n1} & 2.2B (1.34B in the VLM) & VLM: 10\,Hz; action
head: 120\,Hz; 63.9\,ms per 16-action chunk & Timed on an L40 data-centre
GPU & Real, human-video and synthetic \\
$\pi_0$~\cite{pi0} & 3.3B & Up to 50\,Hz & Mobile robot: off-board over
Wi-Fi & About 10{,}000\,h demonstrations \\
Gemini Robotics On-Device~\cite{gemini-ondevice} & Not disclosed &
``Low-latency'', no figure & On the robot; works with no network &
Adapts from 50--100 demos \\
HuGe~\cite{huge-2026} & Qwen-based VLM; total not stated & 50\,Hz action
timeline & Not stated & Over 2{,}000\,h egocentric video \\
GPT-6 Astra~\cite{astranomy-2026} & Not disclosed & About 18.6\,s per call &
Cloud API & No robot-specific training \\
\bottomrule
\end{tabularx}
\caption{Size, speed and location of the reviewed robot brains, as stated in
each source.}
\label{tab:sizes}
\end{table}


\noindent
The trade between model size and reaction speed is regular.  Helix~02's
10M-parameter System~0 runs at 1\,kHz, Helix's 80M-parameter System~1 at
200\,Hz, the 2--7B models at about 7--10\,Hz, and the cloud frontier model
Astra takes about 18.6\,s per decision.  Each step up in size buys more
reasoning and costs about an order of magnitude in speed.

The latency budget then decides which layer a model can be.  Anything that
must react within about 100\,ms has to be a small network on the robot.  Only
task planning can tolerate cloud delay, which is why every Astra system in
Section~\ref{sec:astra} used Astra as the planner only.

An offline requirement rules out cloud models completely, so the top layer
must be a model of a few billion parameters that fits on the robot's computer,
as in Helix's onboard 7B System~2 or Gemini Robotics On-Device, which is
designed to work with intermittent or zero connectivity.  For scale, NVIDIA's
Jetson AGX Thor, a computer built for humanoids, gives up to 2{,}070 TFLOPS
(sparse FP4) within a 40--130\,W power budget and 128\,GB of
memory~\cite{jetson-thor}.  That is enough for several-billion-parameter
models, but it uses battery power that would otherwise move the robot.

The split also gives a second argument for the hierarchical stack: because the
layers are separate, the fast layers can stay on the robot while the slow
planner runs in the cloud when a network is available.  I would also put a
smaller onboard planner in the design as a fallback, so the robot keeps
working, more simply, when the connection drops.

%-----------------------------------------------------------------
\section{The best solution and why}

In my view the best current solution is the \textbf{hierarchical learned
stack}: a slow reasoning model (System~2), a fast visuomotor policy
(System~1) and a learned whole-body controller (System~0), trained so that the
layers share a common interface.  The most complete published example is
\textbf{Figure's Helix~02 on Figure~03}.  For a factory buyer today,
\textbf{Atlas} is the most deployment-ready body, and its pairing with
DeepMind's models points to the same architecture.  I choose this solution for
four reasons.

It matches the physics of the problem.  Balance must be corrected within
milliseconds, grasping within tens of milliseconds, and task planning within
seconds, and one network running at one rate cannot serve all three.  The
three-tier design gives each timescale its own loop, which is the same
principle as the cascaded inner and outer loops we study in control, applied to
learned components.  Section~\ref{sec:deploy} shows that the same split also
decides where each layer can run: the fast loops on the robot, the slow planner
on the robot or in the cloud.

It joins walking and manipulation.  Older humanoids switched between separate
locomotion and arm controllers, so they walked, stopped and only then reached.
A learned System~0 lets the robot walk while carrying and balance during
contact; the dishwasher demonstration is the strongest public evidence of this.

It scales with data rather than engineering.  A new task needs demonstrations
instead of a new controller, and GR00T~N1 and HuGe show that human video and
synthetic data can supplement scarce robot data.

Independent groups converged on it.  GR00T~N1, Gemini Robotics~1.5, HuGe and
even the GPT-6 Astra studies in Section~\ref{sec:astra} all place a slow
planner above faster controllers.

Helix~02 therefore ranks first as an integrated system, with a clear
qualification: its evidence is company demonstrations, with no public benchmark
or paper.  If openness and reproducibility are weighted more heavily,
GR00T~N1 (open weights) and Gemini Robotics~1.5 (detailed technical report)
are better references for how the architecture works.

%-----------------------------------------------------------------
\section{Challenges in improving humanoid robots}

LLMs and AI agents are already part of everyday software work, writing code
and answering customers.  Humanoids run the same kind of model as their
``brain'' and are still mostly at the pilot stage.  The public record is
small.  Agility Robotics and the logistics company GXO signed what they called
the first humanoid Robot-as-a-Service contract in 2024~\cite{agility-wiki}.
Figure reports about 740 robots operating by the end of June 2026, which is a
company figure~\cite{figure-wiki}.  Tesla is still installing its first
Optimus production lines, and analysts expect external sales no earlier than
the second half of 2027~\cite{investorplace-optimus}.

Four differences between software agents and robots account for most of the
gap.  A wrong line of code can be reviewed, tested and undone, while a wrong
physical action can break equipment or hurt someone and cannot be undone.  A
software agent can be copied to a million users almost for free, while every
robot is hardware that must be built, charged, maintained and insured.  A
humanoid also has to beat a working alternative, a fixed automation cell or a
human worker, on total cost and reliability, and both already work well.
Finally, robot data is scarce, and there is no safety standard yet for
certifying the robot.

LLM progress does reach robots quickly, but only at the top of the stack, as
the planner.  The body, the fast control loops and safety are what still limit
commercial use.

Data is scarce and the embodiment gap is wide.  Language models learned from
trillions of words, while humanoid policies have thousands of hours of
demonstrations.  Teleoperation is slow and expensive, and human video does not
contain the robot's joint torques or contact forces; HuGe's need for simulation
replay and arm re-rendering shows how much work is required just to make human
data usable.

Reliability over a long shift is a second problem.  A demo succeeding once is
very different from 99.9\% success over an eight-hour shift.  Learned policies
can fail in surprising ways when lighting, objects or clutter differ from
training data, and those failures are hard to predict in advance.

Safety and certification are unsettled.  A humanoid that falls when power is
lost is a new hazard class, and existing standards do not cover it: ISO~10218
does not address mobility, ISO~3691-4 excludes these robots, and ISO~13482 is
not for industrial sites.  The dedicated standard for dynamically stable
industrial mobile robots, ISO~25785-1, is still in
development~\cite{iso-25785}.  It is also unclear how to certify a neural
network controller whose behaviour cannot be fully specified in advance, and
general models make this harder.  In Robocurve's RoboHarm test (September
2026), Claude Fable~5.1, GPT-6 Astra and the robot model MolmoAct2 controlled
real arms and were each given 5 instructions a safe robot should refuse, 20
times each~\cite{roboharm}.  Fable refused 20 of 100, all on one task (stabbing
a doll); Astra refused 2; MolmoAct2 refused none.  Fable and Astra each pushed
a metal screwdriver into a toaster 6 and 7 times out of 20.  The chatbot safety
training of these models does not carry over to physical actions, so a separate
safety layer on the robot is still needed.

Dexterous hands remain fragile and costly, and tactile sensing is only now
entering learned policies.  Onboard compute also competes with the battery used
for motion (Section~\ref{sec:deploy}), and unit cost must fall far enough to
beat an automation cell or a human worker.  Evaluation is the last gap: most
impressive results are company videos, and without shared benchmarks and
independent tests, systems are hard to compare.

%-----------------------------------------------------------------
\section{Future directions}

Data collection comes first: egocentric human video (as in HuGe), synthetic
data and world models should reduce the dependence on slow teleoperation.
Fleets connected through software such as Orbit can pool corrections, so one
robot's lesson improves all of them.  Learned policies wrapped by verifiable
safety filters (balance or collision constraints, as FAST's centre-of-mass
objectives hint) are the likely route to certification.  Completing
ISO~25785-1 and shared benchmarks such as HumanoidArena would move the field
from demos to measurable progress, and distilled VLAs on dedicated
accelerators will let more of the stack run offline within the latency budgets
of Section~\ref{sec:deploy}.

%-----------------------------------------------------------------
\section{What I learned from this review}

The course material is still in use, in new places.  I expected AI to have
replaced classical robotics; it has not.  The kinematics we did in class
(frames, D--H parameters, forward kinematics) is what retargeting uses to map a
human motion onto the robot's joints, and dynamics and balance quantities such
as the centre of mass are fed into the learned whole-body controller.
Estimating where the robot's floating base is and how fast it moves is still a
filtering problem, like the Kalman filter in Part~2.

Timescale is the main design choice.  Every strong system I read about splits
the work into a slow planner, a middle layer and a fast controller, which is
the cascaded-loop idea from control, only now each loop is a neural network.

Data turned out to be the bottleneck.  Many recent papers look like algorithm
papers, but what they do is find new data: human video, simulation, or a
cheaper way to teleoperate.

Not all sources are equal.  I learned to check what kind of evidence I was
reading.  A peer-reviewed paper, a preprint, a project page and a company video
support a claim to very different degrees, and most humanoid ``results'' are in
the last two groups.

Two things surprised me.  First, GPT-6 Astra beat a model trained on robot data
(74.8\% against 28.0\% success in simulation) without any robot training at
all, yet one call took about 18.6~s; the smartest part of the system is also
the slowest, so it can only work with fast controllers below it.  Second, in
the RoboHarm test the general models refused very few unsafe physical commands,
even though they refuse harmful requests in chat.  I had assumed their safety
training would carry over to the robot, and it does not.

%-----------------------------------------------------------------
\section{Effect of this assignment on my understanding and interest}

My understanding of robotics and intelligent systems did increase.  In the
lectures, kinematics, sensors and filtering felt like separate topics, and
this review showed me how they connect.  Moving an arm to a certain place is
described by frames, transforms and joint angles, and what the sensors perceive
is turned into the same kind of description: a camera image becomes an object
pose, and noisy measurements are fused into a state estimate.  These
expressions are the structured data that a robot ``brain'' learns from and
produces.  A VLA model is trained on camera images paired with joint or
end-effector targets, and its output is again a target pose or joint command
that the lower-level controller can follow.  Learning has not replaced the
mathematics from the course; it is the common language between perception, the
brain and the body.  Without it, there would be nothing well defined for the
model to train on or to output.

My interest in a career in robotics and intelligent systems also increased.  I
have always wanted to work in AI infrastructure, the systems that make large
models possible to train and run, and before this assignment I saw that mainly
as a software and language-model topic.  This review showed me that embodied
intelligence needs the same work at an even larger scale.  The numbers in
Sections~3 and~4 make it concrete: one GPT-6 Astra call takes about 18.6~s,
one agent episode uses about 30.9 million tokens (about US\$36), and the fast
controllers must fit on an onboard computer with a limited power budget, while
robot data is scarce and expensive to collect.  A humanoid improves in a
cycle: train the model, deploy it, collect new data from what the robot does,
especially its failures, and retrain.  Every step of that cycle depends on
infrastructure: training clusters, data pipelines for video and sensor logs,
simulation at scale, and model compression so that inference is fast enough on
the robot.  Optimising these steps decides how fast the cycle turns, and so how
fast humanoids improve.  That gives me a clear way to combine my interest in AI
infrastructure with robotics, and it is the direction I would like to pursue.

%-----------------------------------------------------------------

\begin{thebibliography}{99}

\bibitem{bd-atlas-2026}
Boston Dynamics, ``Boston Dynamics unveils new Atlas robot to revolutionize
industry,'' company announcement, 5 Jan.\ 2026.
\url{https://bostondynamics.com/blog/boston-dynamics-unveils-new-atlas-robot-to-revolutionize-industry/}

\bibitem{helix-02}
Figure AI, ``Introducing Helix 02: Full-body autonomy,'' company
announcement, Jan.\ 2026.
\url{https://figure.ai/news/helix-02}

\bibitem{fast-2026}
Z.~Wang \emph{et al.}, ``General humanoid whole-body control via pretraining
and fast adaptation,'' arXiv:2602.11929, Feb.\ 2026.

\bibitem{bfm-survey}
M.~Yuan \emph{et al.}, ``A survey of behavior foundation model:
Next-generation whole-body control system of humanoid robots,''
\emph{IEEE Trans.\ Pattern Anal.\ Mach.\ Intell.}, vol.~48, pp.~4909--4927,
2026, doi:10.1109/TPAMI.2025.3649177 (arXiv:2506.20487).

\bibitem{huge-2026}
Y.~Liu \emph{et al.}, ``HuGe: Scaling humanoid generalist
vision-language-action models with egocentric experience,''
arXiv:2610.34186, Oct.\ 2026.

\bibitem{helix-2025}
Figure AI, ``Helix: A vision-language-action model for generalist humanoid
control,'' company announcement, 20 Feb.\ 2025.
\url{https://www.figure.ai/news/helix}

\bibitem{gemini-ondevice}
Google DeepMind, ``Gemini Robotics On-Device brings AI to local robotic
devices,'' blog post, 24 Jun.\ 2025.
\url{https://deepmind.google/discover/blog/gemini-robotics-on-device-brings-ai-to-local-robotic-devices/}

\bibitem{jetson-thor}
NVIDIA, ``Introducing NVIDIA Jetson Thor, the ultimate platform for physical
AI,'' NVIDIA Technical Blog, 25 Aug.\ 2025.
\url{https://developer.nvidia.com/blog/introducing-nvidia-jetson-thor-the-ultimate-platform-for-physical-ai/}

\bibitem{robocurve-astra}
A.~Menon, S.~Machcha, S.~Zou, T.~K. Chan and J.~Chooi, ``GPT-6 Astra on
robotic manipulation,'' Robocurve, 4 Sep.\ 2026.
\url{https://openai.robocurve.org/gpt-6-astra/}

\bibitem{roboharm}
E.~Sun, S.~Machcha, S.~Zou, T.~K. Chan and J.~Chooi, ``RoboHarm: Do frontier
robot policies refuse unsafe instructions?,'' Robocurve, 18 Sep.\ 2026.
\url{https://robocurve.org/roboharm/}

\bibitem{grok-optimus}
K.~Singh, ``Tesla to integrate xAI's Grok into Optimus, helping bring the
robot to life,'' Not a Tesla App, 9 Jul.\ 2025.
\url{https://www.notateslaapp.com/news/2874/tesla-to-integrate-xais-grok-into-optimus-helping-bring-the-robot-to-life}

\bibitem{investorplace-optimus}
J.~Rose, ``Tesla Optimus, Grok, and TSLA stock: Inside the physical-AI bet
that now defines Tesla,'' InvestorPlace, 7 Sep.\ 2026.
\url{https://investorplace.com/dailylive/2026/09/tesla-optimus-grok-and-tsla-stock-inside-the-physical-ai-bet-that-now-defines-tesla/}

\bibitem{optimus-wiki}
``Optimus (robot),'' Wikipedia, accessed 7 Nov.\ 2026.
\url{https://en.wikipedia.org/wiki/Optimus_(robot)}

\bibitem{figure-wiki}
``Figure AI,'' Wikipedia, accessed 7 Nov.\ 2026.
\url{https://en.wikipedia.org/wiki/Figure_AI}

\bibitem{agility-wiki}
``Agility Robotics,'' Wikipedia, accessed 7 Nov.\ 2026.
\url{https://en.wikipedia.org/wiki/Agility_Robotics}

\bibitem{homebody-2026}
G.~Huh, C.~Gu, T.~E. Truong, C.~K. Liu and G.~Tevet, ``HomeBody,'' project
page, Stanford University and Caltech, Sep.\ 2026.
\url{https://tml.stanford.edu/homebody/}

\bibitem{astranomy-2026}
J.~X. Yang, A.~Khetan and S.~Zhang, ``Astranomy: An empirical study of GPT-6
Astra as a native robot policy,'' NVIDIA Research, CoRL 2026 late-breaking
result, Oct.\ 2026.
\url{https://research.nvidia.com/labs/srl/projects/astranomy/}

\bibitem{sensenova-2026}
SenseNova Joint Research Team \emph{et al.}, ``Has GPT-6-Astra achieved
embodied intelligence?,'' arXiv:2610.22471, Oct.\ 2026.

\bibitem{motionvla-2026}
``MotionVLA: Vision-language-action model for humanoid motion,''
arXiv:2606.15142, Jun.\ 2026.

\bibitem{gemini-15}
Gemini Robotics Team, Google DeepMind, ``Gemini Robotics 1.5: Pushing the
frontier of generalist robots with advanced embodied reasoning, thinking, and
motion transfer,'' arXiv:2510.03342, Oct.\ 2025.

\bibitem{groot-n1}
NVIDIA (J.~Bjorck \emph{et al.}), ``GR00T N1: An open foundation model for
generalist humanoid robots,'' arXiv:2503.14734, Mar.\ 2025.

\bibitem{pi0}
K.~Black \emph{et al.}, ``$\pi_0$: A vision-language-action flow model for
general robot control,'' arXiv:2410.24164, Oct.\ 2024.

\bibitem{figure-03}
Figure AI, ``Introducing Figure 03,'' company announcement, Oct.\ 2025.
\url{https://figure.ai/news/introducing-figure-03}

\bibitem{iso-25785}
ISO, ISO~25785-1, Part~1 of a safety standard for industrial mobile robots
whose stability depends on active control (including legged humanoids),
committee draft, under development.
\url{https://www.iso.org/standard/91469.html}

\end{thebibliography}


\end{document}
